\begin{proof}
\begin{enumerate}[label={$\langle1\rangle$\arabic*.},leftmargin=*]
\item 
\noindent\textbf{Step 1.} In coordinates with $S=(z_1\cdots z_r=0)$, $1\le r\le3$, the map sends a vector field $\sum a_i\partial_{z_i}$ to the normal coefficient $a_i|_{z_i=0}$ on each branch $i\le r$. The separate coefficients lift independently to holomorphic functions, making the map surjective. Its kernel consists precisely of fields with $z_i\mid a_i$ for $i\le r$, which is $T_Y(-\log S)$. Coordinate changes act on each coefficient through the normal-line transition and permute the branches; hence the map and its exactness descend canonically to the normalization. The immersion tangent sequence on $\widetilde S$ gives $\det\nu^*T_Y=\det T_{\widetilde S}\otimes N_\nu$, proving the adjunction expression. Tensoring by the locally free line bundle $L$ preserves exactness. Finite pushforward has no higher direct images, so the long exact sequence in coherent cohomology gives \eqref{cdp:eq:C.4}. The smooth compact surface $\widetilde S$ has the usual Serre pairing, which gives the stated dual target. In particular the vanishing of $H^3(Y,\mathcal T\otimes L)$ makes this cohomology map surjective; no such vanishing has been suppressed in \eqref{cdp:eq:C.4}.

\end{enumerate}
\end{proof}
